\chapter{Conclusion}
Retailers hold CCTV footage that could answer everyday operational questions but is rarely used for that purpose, and small and mid-sized stores often cannot afford specialised analytics. RetailPulse addresses this with a CPU-friendly, non-biometric, multi-agent design that reuses existing cameras. Research on detection, tracking, and re-identification, and a feasibility triage, shaped the scope: a single-store, three-camera prototype focused on person detection, role identification and pose-checked interaction detection.

At the mid-report stage the requirements and architecture are defined. The web application with manager and employee roles, secure session handling and person-analysis prototypes form the current implementation; report export remains a planned feature on the tested branch. An interaction pipeline combines the trained employee/customer detector, body-height-normalised proximity and a YOLO pose judge that confirms an interaction only when someone faces the other person; on the one labelled clip the judge raises interaction precision from 0.50 to 1.00, and the detector reaches mAP50 0.969 on 455 held-out images. A floor-distance CLI passed seventeen focused checks before it was descoped, and 44 current interaction unit tests pass. Dashboard vision metrics still rely on prototype data. The main risks are CPU throughput, classification accuracy, and interaction interpretation. FYP-2 will integrate the vision agents, evaluate them on labelled store footage, and harden the system for the final demonstration.
